%!TeX root=../Memory.tex


\chapter{Estat de l'art}
\label{ch:estat-art}

La votació electrònica és un camp amb una història llarga i un nombre significatiu d'experiències reals desplegades a
diferents països i contextos institucionals.
Aquest capítol revisa, en primer lloc, els casos d'ús desplegats a escala nacional que han marcat el debat internacional
sobre la viabilitat del vot electrònic; en segon lloc, els sistemes acadèmics que defineixen el referent criptogràfic
actual de la votació verificable; i, en tercer lloc, els projectes contemporanis que estan combinant criptografia
avançada amb infraestructura \emph{blockchain}.
L'objectiu és situar el projecte \emph{Demochain} en un panorama tècnic concret abans d'entrar en l'anàlisi i el disseny
del sistema.

%-----------------------------------------------------------------------


\section{Experiències nacionals de votació electrònica}
\label{sec:experiencies-nacionals}

La transformació digital de processos col·lectius --consultes, primàries, eleccions internes d'entitats, referèndums--
ha posat de manifest una tensió de fons: els sistemes de votació electrònica desplegats fins ara tendeixen a substituir
un punt únic de confiança (el responsable del recompte físic) per un altre (el responsable del sistema informàtic).
El problema no és intrínsec a la digitalització, sinó al model arquitectònic.
Un servidor central que rep, emmagatzema i recompta vots és, per definició, una entitat en qui s'han de dipositar
fiançament tant la integritat del recompte com la privacitat del vot individual.

L'experiència internacional dels darrers vint anys mostra el ventall de resultats que ha donat aquest model.

\subsection{Estònia: l'e-voting com a cas paradigmàtic}
\label{subsec:estonia:-l'e-voting-com-a-cas-paradigmatic}

Estònia és el cas paradigmàtic d'adopció a gran escala.
El sistema \emph{i-Voting} permet el vot a Internet per a eleccions nacionals vinculants des de
2007~\cite{e-estonia2024voting} i la proporció de votants per Internet ha crescut de manera sostinguda fins a superar el
50\% en eleccions recents.
El sistema es recolza en el document d'identitat digital nacional (eID) i en un protocol de doble sobre electrònic: el
vot intern es xifra amb la clau pública de la comissió electoral i s'envolta d'una signatura amb la clau privada del
votant.

Tot i ser un cas d'èxit pel que fa al desplegament massiu, l'e-voting estonià ha estat objecte d'anàlisi de seguretat
externa.
Springall et al. \cite{springall2014estonia} van publicar el 2014 un estudi detallat que identifica vulnerabilitats
serioses en aspectes operacionals i procedimentals --la gestió del \textit{software} del client, la construcció de les imatges
del sistema operatiu d'arrencada, els controls físics al centre d'operacions-- que il·lustren la dificultat de garantir
totes les propietats desitjades en un sistema centralitzat.
Les autoritats estonianes van respondre amb millores incrementals, però el model arquitectònic centralitzat es manté.

\subsection{Suïssa: el cas SwissPost i la suspensió del 2019}
\label{subsec:suissa:-el-cas-swisspost-i-la-suspensio-del-2019}

Suïssa va desplegar entre 2018 i 2019 el sistema de \emph{e-voting} de SwissPost, basat en una implementació de Scytl.
El sistema declarava oferir verificabilitat universal mitjançant proves de coneixement zero sobre les decisions de
mescla (\emph{mixnet}) que reordenen els vots.

El març de 2019 un equip de criptògrafs independents~\cite{lewis2019swisspost} va publicar una anàlisi que demostrava
que les proves de Bayer-Groth implementades empraven \emph{commitments} amb \emph{trapdoor}, una propietat que en
principi hauria de ser absent.
Aquesta debilitat permetria, en escenaris específics, alterar el resultat sense que la verificació detectés cap
incidència.
La publicació va portar a la suspensió del sistema.
Només va tornar a desplegar-se anys més tard amb una reescriptura completa i un règim d'auditoria pública del codi.

\subsection{Estats Units: l'aplicació Voatz a West Virginia}
\label{subsec:estats-units:-l'aplicacio-voatz-a-west-virginia}

Als Estats Units, l'aplicació \emph{Voatz} es va emprar en eleccions federals primàries de l'estat de West Virginia el
2018, dirigida inicialment a votants residents a l'estranger.
Voatz promocionava l'ús de tecnologia \emph{blockchain} per garantir la integritat dels vots, però l'arquitectura
interna era essencialment centralitzada: la \emph{blockchain} s'emprava com a capa de registre d'auditoria secundària,
no com a capa primària d'execució.

L'anàlisi de seguretat publicada per Specter, Koppel i Weitzner del MIT el 2020~\cite{specter2020voatz} va revelar
vulnerabilitats que permetien alterar, observar i invalidar el vot dels usuaris, així com problemes estructurals en el
model de seguretat declarat per l'aplicació.
La conclusió dels autors --reflectida al títol de l'article, \emph{The Ballot is Busted Before the Blockchain}-- és que
la mera incorporació de tecnologia \emph{blockchain} no garanteix la integritat del sistema si l'arquitectura de
l'aplicació presenta defectes en altres capes.

\subsection{Lliçó general}\label{subsec:llico-general}

Aquests tres casos no qüestionen la viabilitat de la votació electrònica en si: el problema no és la digitalització,
sinó com es distribueix la confiança.
La pregunta de fons és si és possible construir un sistema en què cap entitat individual --ni l'administrador, ni el
desenvolupador, ni el proveïdor d'infraestructura-- pugui alterar o manipular el resultat sense que això sigui
detectable.
És precisament aquesta la promesa de la tecnologia \emph{blockchain} aplicada a la votació, sempre que s'apliqui amb
honestedat i no com a element de \emph{marketing}.

%-----------------------------------------------------------------------


\section{Sistemes acadèmics de votació verificable}
\label{sec:sistemes-academics}

L'aplicació de la criptografia a la votació electrònica verificable té una tradició acadèmica madura.
Els projectes següents han estat referents directes o indirectes durant la fase de recerca del nostre projecte.

\subsection{Helios}
\label{subsec:helios}

Helios~\cite{adida2008helios}, presentat per Ben Adida al USENIX Security Symposium de 2008, és un sistema obert de
votació web amb verificabilitat universal.
La seva contribució fonamental és la combinació de \emph{mixnets} per garantir l'anonimat dels vots i xifratge ElGamal
amb proves de coneixement zero per permetre que qualsevol persona pugui verificar que els vots s'han inclòs al recompte
sense conèixer-ne el contingut.

Helios s'ha emprat per a eleccions internes d'institucions acadèmiques: la \emph{International Association for
Cryptologic Research} (IACR), la Universitat Catòlica de Lovaina, la Universitat de Princeton, entre d'altres.
El seu model d'amenaces, però, es basa en un servidor central que executa el \emph{tally} (el càlcul del recompte): una
sola entitat pot bloquejar el procés si es nega a publicar el resultat.
Aquest punt feble va motivar el desenvolupament de Belenios.

\subsection{Belenios}
\label{subsec:belenios}

Belenios~\cite{cortier2019belenios}, desenvolupat per l'equip de Véronique Cortier a INRIA Nancy, és una evolució
conceptual d'Helios.
Aporta dues novetats principals: xifratge \emph{threshold} ElGamal i una cerimònia de desxifrat distribuïda entre
múltiples \emph{trustees}.
Amb aquest esquema, cap \emph{trustee} per separat pot desxifrar un vot individual; només una majoria qualificada
(\emph{p.ex.}, 3 de 5) pot reconstruir la clau privada per descomprimir el recompte agregat.

Aquesta arquitectura elimina el punt únic de fallada del \emph{tally authority} d'Helios, però introdueix una
complexitat operativa addicional: la cerimònia de desxifrat requereix coordinació entre els \emph{trustees}, cadascun
amb la seva clau, i si massa \emph{trustees} es perden o es neguen a participar, el resultat no es pot calcular.
Belenios s'empra activament per a eleccions universitàries franceses.

\subsection{Microsoft ElectionGuard}
\label{subsec:microsoft-electionguard}

ElectionGuard~\cite{microsoft2019electionguard} és un SDK obert publicat per Microsoft Research el 2019 per a votació
verificable \emph{end-to-end}.
La seva proposta principal és aportar una capa criptogràfica reutilitzable que es pot integrar amb sistemes existents
(urnes electròniques, sistemes de votació per correu) sense substituir-los.
Empra, com Belenios, ElGamal \emph{threshold} amb proves de coneixement zero.

ElectionGuard s'ha desplegat a un nombre creixent d'eleccions reals als Estats Units en mode pilot, sempre com a capa de
verificabilitat adjunta a un sistema de votació primari.
La seva filosofia és deliberadament conservadora: no pretén substituir els sistemes existents, sinó dotar-los d'una capa
de verificabilitat criptogràfica.

%-----------------------------------------------------------------------


\section{Projectes contemporanis sobre blockchain}
\label{sec:projectes-blockchain}

Helios, Belenios i ElectionGuard han assolit un grau alt de maduresa criptogràfica, però mantenen un component
centralitzat clar: l'autoritat electoral és sempre una entitat coneguda i en última instància responsable de l'execució
del sistema.
La novetat conceptual que aporten les xarxes \emph{blockchain} és la possibilitat de substituir aquesta autoritat per un
protocol distribuït de consens, on un conjunt de nodes independents executen i verifiquen la lògica electoral sense que
cap d'ells pugui modificar-la unilateralment.

\subsection{DAVINCI Protocol (Vocdoni)}
\label{subsec:davinci-protocol-(vocdoni)}

El projecte més rellevant en aquesta línia, en el moment de redactar aquesta memòria, és \emph{DAVINCI
Protocol}~\cite{vocdoni2026davinci}, desenvolupat per Vocdoni.
Vocdoni és un equip amb experiència contrastada des de 2018 en governança descentralitzada, que ha treballat amb partits
polítics, governs locals, entitats cíviques i DAO\@.

DAVINCI combina una capa \emph{Layer 2} sobre Ethereum amb una xarxa descentralitzada de \emph{sequencers}, una màquina
d'estat verificada amb proves \acp{ZK-SNARK} i xifratge \emph{threshold} ElGamal per garantir privacitat del vot.
Inclou també mecanismes de resistència a la coacció (\emph{receipt-freeness}) que impedeixen al votant demostrar a un
tercer que ha votat, una propietat fonamental per prevenir la compra de vots i la coerció en entorns no controlats.

El projecte es publica sota llicència AGPL-3.0 i és, dels sistemes revisats, el més proper conceptualment al que es
voldria arribar a llarg termini amb \emph{Demochain}: materialitza, en producció, bona part de les propietats
criptogràfiques que s'han hagut de deixar fora de l'abast del nostre \ac{MVP} (vegeu secció~\ref{sec:abast} del
capítol~\ref{ch:introduccio}).

\subsection{Lliçons recollides}
\label{subsec:llicons-recollides}

La revisió del panorama actual permet identificar tres tendències clares en l'estat de l'art:

\begin{enumerate}
    \item Els sistemes acadèmics consolidats (Helios, Belenios, ElectionGuard) aporten primitives criptogràfiques
    sòlides, però mantenen un component centralitzat operatiu.

    \item Els projectes contemporanis sobre \emph{blockchain} (DAVINCI) combinen aquestes primitives amb infraestructura
    distribuïda, però requereixen un disseny d'infraestructura significativament més complex (xarxes pròpies de
    \emph{sequencers}, \emph{zkVMs} dedicats).

    \item Els intents d'adopció massiva amb tecnologia comercial sense rigor criptogràfic (Voatz, casos primerencs de
    SwissPost) han tendit a fracassar quan han estat sotmesos a auditoria externa.
\end{enumerate}

El nostre projecte se situa, deliberadament, en una posició intermèdia: aprofita una infraestructura \emph{blockchain}
existent i pública (Algorand) com a capa de votació, en lloc de construir-ne una de pròpia, i empra una segona
\emph{blockchain} pública (Ethereum) com a capa de notaria independent.
La capa criptogràfica avançada queda explícitament identificada com a treball futur, sense ocultar-ho en cap moment ni
reclamar propietats que el sistema actual no garanteix.
La justificació detallada d'aquesta tria arquitectònica es recull al capítol de disseny.
